\chapter{Dinámica, estabilidad y límites termoestructurales}\label{ch:13}
\objetivos{Linealizar un equilibrio concentrado; separar estabilidad de presión media y combustión acústica; relacionar escalas térmicas y mecánicas sin utilizar modelos elementales como certificación estructural.}
\section{Estabilidad local de un equilibrio concentrado}
En el modelo normalizado \eqref{eq:solidNorm}, tomemos área fija $a=1$ y equilibrio $P=1$. Escribamos $P=1+\delta$, con $|\delta|\ll1$. La expansión binomial da
\[
 P^n=1+n\delta+O(\delta^2),
\]
y la ecuación linealizada es
\[
 \frac{d\delta}{d\theta}=(n-1)\delta.
\]
Por integración,
\begin{equation}\label{eq:solidStability}
 \delta(\theta)=\delta(0)e^{-(1-n)\theta}.
\end{equation}
Para $n<1$ las perturbaciones pequeñas decaen dentro de este modelo de un único estado, volumen fijo, temperatura constante y descarga ahogada. El tiempo de relajación dimensional es $\tau_0/(1-n)$. Acercarse a uno aumenta a la vez sensibilidad estática y lentitud de recuperación.

Para $n=1$, la linealización es neutra y la existencia del equilibrio depende del coeficiente de generación respecto a descarga. Para $n>1$, la perturbación crece en el modelo. Estas conclusiones no autorizan a clasificar un motor real como estable o inestable sólo por su exponente medio. La dinámica térmica de superficie, el acoplamiento acústico y los retardos no están representados.
\section{Respuesta a una perturbación geométrica}
Si $a=1+\epsilon$ y $P=1+\delta$, a primer orden,
\[
 \frac{d\delta}{d\theta}=\epsilon-(1-n)\delta.
\]
Para escalón constante y condición $\delta(0)=0$,
\[
 \delta(\theta)=\frac{\epsilon}{1-n}
 \left[1-e^{-(1-n)\theta}\right].
\]
El límite $\epsilon/(1-n)$ coincide con la expansión lineal de la solución cuasiestacionaria $(1+\epsilon)^{1/(1-n)}-1$. Esa coincidencia es una verificación cruzada entre cálculo diferencial y equilibrio algebraico. Para perturbaciones grandes se compara con la ecuación no lineal, porque los términos de orden superior pueden ser relevantes.
\section{Presión media y modos acústicos}
Una cámara real admite distribuciones espaciales. Las fluctuaciones de presión pueden acoplarse a liberación de calor y a la respuesta de la superficie. Un modelo acústico elemental sin flujo medio satisface
\[
 \frac{\partial^2p'}{\partial t^2}=a^2\nabla^2p'
\]
con condiciones de contorno apropiadas. Las frecuencias no dependen sólo del volumen: también de la geometría, impedancias, velocidad del sonido y acoplamientos. Un gas caliente no uniforme requiere una descripción más rica.

La idea energética es que la liberación de calor en fase con una presión elevada puede alimentar una oscilación. Para una descripción termoacústica, el signo de un promedio espacial y temporal de $p'\dot q'$ es un indicador de aporte o extracción de energía acústica, bajo las hipótesis que permiten construir ese balance. No es un criterio completo si se omiten pérdidas, fronteras y transporte. Este razonamiento se retoma para motores líquidos.\cite{MIT22,MIT23}
\section{Flujo erosivo y partículas}
Una corriente gaseosa sobre una superficie puede modificar su intercambio térmico y la tasa de regresión. Esa dependencia no está contenida en una ley que sólo utiliza presión. Sustituirla por un cambio arbitrario de $n$ puede ajustar una curva, pero no identifica ni predice el mecanismo fuera de los datos.

Los productos condensados pueden intercambiar momento y calor con el gas y acumularse en ciertas regiones. Las partículas con velocidad inferior a la del gas reducen el momento axial respecto a un modelo de equilibrio mecánico perfecto. La masa total sigue conservándose: la diferencia está en cómo se reparte su energía y velocidad. Un modelo bifásico requiere tamaños, propiedades y escalas de relajación; no se deduce con una simple corrección de masa molar.\cite{MIT06}
\section{Integridad y escalas térmicas}
La envolvente soporta carga de presión y restricciones térmicas. Como referencia mecánica elemental, una membrana cilíndrica delgada, lejos de extremos y discontinuidades, cumple
\[
 \sigma_\theta\simeq\frac{pr}{t},\qquad
 \sigma_z\simeq\frac{pr}{2t}
\]
para extremos cerrados y carga interna diferencial $p$. Estas igualdades se obtienen equilibrando la fuerza de presión con las tracciones de membrana. Exigen $t/r\ll1$, geometría regular y ausencia de concentraciones. No representan juntas, cierres, gradientes térmicos ni materiales compuestos anisótropos.

La conducción transitoria tiene profundidad característica $\ell_{th}\sim\sqrt{\alpha t_{exp}}$. Un espesor grande respecto a esa escala puede mantener una zona posterior relativamente fría durante un intervalo, pero no establece por sí solo una temperatura admisible. Hace falta resolver condiciones iniciales y de contorno, propiedades dependientes de temperatura y posibles cambios de fase.

La ablación consume o transforma material y modifica la frontera térmica. Una protección aislante reduce flujo hacia la estructura; una masa absorbente acumula energía sensible; una superficie radiante rechaza energía al entorno. Sus balances son diferentes. Hablar genéricamente de refrigeración sin distinguir el destino de la energía impide comprobar conservación.\cite{MIT10}
\section{Integración de prestaciones y dispersión}
El impulso total se calcula como $I=\int C_F(t)p_c(t)A_t(t)\,dt$. Una curva de área progresiva no permite integrar el empuje sin la evolución temporal de presión, garganta y ambiente. Si se emplea un promedio, debe declararse qué señal se promedia y durante qué intervalo.

La dispersión de propiedades puede amplificarse por la sensibilidad $1/(1-n)$. En un estudio docente se pueden propagar incertidumbres mediante derivadas o muestreo de parámetros normalizados, conservando sus correlaciones y dominios. Un conjunto de curvas numéricas no representa una distribución real de fabricación si los parámetros no provienen de datos.
\figura{c13}{Estabilidad local de un único estado con temperatura y volumen fijos. No evalúa modos acústicos ni certifica estabilidad de combustión.}
\section{Caso desarrollado}
\begin{studiumworked}{M13. Respuesta lineal a un pequeño cambio de área · Intermedio}
\textbf{Enunciado.} El modelo normalizado $dP/d\theta=aP^n-P$ tiene $n=0.4$, $P(0)=1$. El área salta a $a=1.02$. Obtenga la respuesta lineal en $\theta=3$, el equilibrio lineal y el equilibrio no lineal.

\textbf{Datos.} Exponente: $0.4$; Escalón: $0.02$; Tiempo relativo: $3$.

\textbf{Incógnitas.} Magnitudes solicitadas y cierre físico.

\textbf{Hipótesis.} Modelo académico con hipótesis del capítulo. Datos originales; no representan hardware ni ensayo.

\textbf{Modelo físico.} Conservación y relaciones constitutivas declaradas.
\[P=1+\delta,\quad d\delta/d\theta=\epsilon-(1-n)\delta\]
\[\delta(\theta)=\epsilon[1-e^{-(1-n)\theta}]/(1-n)\]
\[P_{eq}=a^{1/(1-n)}\]
\textbf{Desarrollo.}
\begin{enumerate}
\item Linealizar alrededor de uno elimina productos de perturbaciones. El equilibrio lineal es $\delta_{eq}=0.02/0.6=0.03333333$.
\item A tiempo relativo tres, $\delta=0.02782337$ y $P\simeq1.027823$. El tiempo de relajación es $1/(1-n)=1.666667$ unidades de $\theta$.
\item El equilibrio no lineal tiene $P_{eq}-1=1.02^{1/0.6}-1=0.03355506$. La diferencia con $0.03333333$ es de orden cuadrático en el escalón.
\end{enumerate}
\textbf{Operaciones para contrastar el cálculo.}
\begin{enumerate}

\item Magnitud 1.

\[\fitmath{y_{1}=\frac{0.02}{0.6}\cdot \left(1-\exp\left(-0.6\cdot 3\right)\right)\approx 0.02782337\,}\]

\item Magnitud 2.

\[\fitmath{y_{2}={1.02}^{\frac{1}{0.6}}\approx 1.033555\,}\]

\end{enumerate}

\textbf{Resultados.} $P(3)\simeq1.027823$; $P_{eq,lineal}=1.033333$; $P_{eq,no\ lineal}=1.033555$.

\textbf{Interpretación.} El resultado corresponde a las hipótesis indicadas; su interpretación requiere conservar fronteras y unidades.

\textbf{Comprobación.} La respuesta arranca con delta cero y tiende al equilibrio lineal; el autovalor -0.6 es negativo.

\textbf{Límites.} Estabilidad local del modelo concentrado, no estabilidad acústica de un motor.
\end{studiumworked}

\section{Ejercicios propuestos}
\begin{enumerate}
\item Derive el tiempo de relajación y compare $n=0.30$ con $n=0.80$ en unidades de $\tau_0$.
\item Calcule el error de la aproximación lineal del equilibrio para $a=1.1$, $n=0.4$.
\item Explique por qué una presión media que vuelve al equilibrio no descarta una oscilación acústica creciente.
\item Distinga aislamiento, acumulación sensible y ablación escribiendo un término energético característico de cada mecanismo.
\end{enumerate}
\resumen{La estabilidad de un estado concentrado es una conclusión limitada a sus variables e hipótesis. La cámara puede tener modos espaciales ausentes del modelo medio. Integridad, transferencia térmica y dispersión requieren información adicional; las relaciones elementales permiten interpretar tendencias y comprobar balances, no certificar un motor.}
